% Technical note for collaborators: retrieval-conditioned graded dispatch.
% Companion to redundancy_note.tex. ICLR 2027 template, no title block.
% Compile with pdflatex.
\documentclass{article}
\usepackage{iclr2027_conference,times}
\iclrfinalcopy
\usepackage{amsmath,amssymb}
\usepackage{enumitem}

\begin{document}

\paragraph{Problem.}
On a cache hit, the executor must decide how much of the action head's
denoising schedule to skip. Let $N$ be the number of denoising steps of the
flow-matching action expert and let $\tau \in \{0, 1, \dots, N\}$ denote the
skip depth: the integration is initialized at the retrieved intermediate
state of depth $\tau$ and runs the remaining $N - \tau$ steps. In this view
every hit is a warm start of some depth, and the familiar tiers are its
extreme cases: full replay is the special case $\tau = N$, in which the
retrieved chunk is executed directly, and full inference is the case
$\tau = 0$, generation from pure noise. Convention: $\tau$ counts
\emph{skipped} steps; the executor's start parameter is the remaining
integration fraction, $\mathtt{start\_t} = (N - \tau)/N$. We fix this mapping
here because sign errors between data generation, plots, and the executor are
otherwise easy to make. The dispatch question is therefore how deeply to skip
on each hit. Existing rules are coarse: our current dispatch quantizes $\tau$
to three values, while fixed-depth rules apply one global $\tau$ to every hit.
We instead seek a per-step rule that allocates the remaining denoising
computation without training a separate policy.

\paragraph{Two statistics.}
Both inputs are byproducts of the retrieval already performed. The first is
the fused similarity $s_t$, which measures whether the current state is
covered by the library. The second measures whether the retrieved
neighborhood agrees on what to do. Let $a_{(1)}, \dots, a_{(k)}$ be the
action chunks of the top-$k$ retrieved entries, $\bar{a}$ their mean, and $W$
a fixed diagonal (or Mahalanobis) weighting that standardizes the action
dimensions; define the local action disagreement
\begin{equation}
v_t \;=\; \frac{1}{k}\sum_{j=1}^{k} \bigl\| W\bigl(a_{(j)} - \bar{a}\bigr) \bigr\|^2 .
\label{eq:v}
\end{equation}
Because the top-$k$ entries come from neighboring states rather than repeated
samples at the same state, $v_t$ is not a pure measure of aleatoric
uncertainty. It may reflect retrieval error, task-phase aliasing, action
multimodality, teacher stochasticity, and library density. We therefore treat
it as a local disagreement statistic whose predictive value is measured
empirically.

\paragraph{Random loss under coupled noise.}
Let $c$ denote the current conditioning (observation and prompt), let
$x_{j,\tau}$ be the retrieved intermediate of depth $\tau$ from library entry
$j$, and let $z_j$ be the initial noise of the cached trajectory that
produced it. Write $F_{c}^{\tau \to N}$ for the probability-flow integration
from depth $\tau$ to completion under conditioning $c$. Define
\begin{equation}
a^{(\tau)} \;=\; F_{c}^{\tau \to N}\bigl( x_{j,\tau} \bigr),
\qquad
a^{(0)} \;=\; F_{c}^{0 \to N}\bigl( z_j \bigr),
\qquad
Y_\tau \;=\; \bigl\| W\bigl(a^{(\tau)} - a^{(0)}\bigr) \bigr\| .
\label{eq:Y}
\end{equation}
The reference $a^{(0)}$ is full generation under the \emph{current}
conditioning from the \emph{same} initial noise $z_j$. The coupling matters
because identical distribution does not mean identical realization: a freshly
sampled $z' \sim \mathcal{N}(0, I)$ may land on a different action branch, and
the measured deviation would then confound the warm-start mismatch with the
teacher's own sampling variability. By construction $Y_0 = 0$: at
$\tau = 0$ the two branches coincide.

The calibration artifact therefore stores, for each cache entry (equivalently,
each action-generation call), either the initial noise tensor or a
deterministic seed. A $50 \times 32$ float32 tensor requires approximately
$6.4$\,KB, which is small relative to the stored denoising intermediates. A
seed is cheaper, but it requires bit-exact regeneration across the software
and hardware configuration used for calibration. Current runtime payloads
need not change: the additional field is required only while constructing the
calibration data. As a preliminary diagnostic, we also measure the teacher's
noise-induced variability through repeated full generations at the same
observation. If teacher--teacher deviation is negligible relative to
warm-start deviation, an uncoupled analysis may be reported separately; its
target is population deviation from a random teacher sample rather than the
paired initialization error defined above.

\paragraph{Conditional quantiles.}
For $\tau \ge 1$ define
\begin{equation}
q_\tau(s, v) \;=\; Q_{1-\alpha}\bigl( Y_\tau \mid s_t = s,\; v_t = v \bigr),
\label{eq:q}
\end{equation}
the conditional quantile of the random loss given the retrieval statistics;
it is a function of $(\tau, s, v)$, not a quantile of an expectation. The
qualitative motivation is empirical rather than theoretical: additional
refinement steps may correct initialization mismatch, and a retrieved branch
of a multimodal target can only be arbitrated by the steps that remain. We
do not derive these effects from smoothness constants of the flow map; a
Lipschitz bound alone does not imply contraction.

\paragraph{Why a threshold is not naive.}
Consider the binary decision ``accept or reject a full replay'' based on a
fixed scalar score $s$, and define the offline-checkable event
\begin{equation}
Z \;=\; \mathbb{1}\{\, Y_N \le \delta \,\},
\label{eq:Z}
\end{equation}
direct replay stays within the prescribed deviation tolerance. If the
likelihood ratio of the conditional densities of $s$ given $Z = 1$ and
$Z = 0$ is monotone in $s$, then by the Neyman--Pearson lemma the threshold
test on $s$ is, at any fixed false-accept rate, optimal among all decision
rules that depend on $s$ alone; equivalently, thresholding traces the ROC
frontier of the score. The statement is conditional on the fixed constraint
and on the information available; it does not assert optimality against rules
with access to more information. We use the deviation event because the
effect of replaying one step on the eventual episode outcome is a closed-loop
counterfactual that offline data cannot label. Assigning the episode outcome
to every constituent step would not resolve this attribution problem. The
surrogate is measurable on logged data, while task success is assessed through
closed-loop evaluation. Improvements over the threshold must therefore come
from one of three sources: additional statistics beyond $s$, better
calibration of $s$, or a graded choice of $\tau$ instead of a binary one. The
construction below addresses all three. The monotone-likelihood-ratio premise
for this surrogate is assessed diagnostically on the fit split (E5).

\paragraph{Monotonicity assumption.}
For $\tau \ge 1$,
\begin{equation}
\text{(A1)}\qquad q_\tau(s, v)\ \text{is nondecreasing in } \tau\
\text{and in } v,\ \text{and nonincreasing in } s .
\label{eq:A1}
\end{equation}
A sufficient and cleaner condition is first-order stochastic dominance of the
conditional laws of $Y_\tau$ along each coordinate. (A1) is an assumption
about the data-generating process, not a theorem; its empirical violation
rate on the fit split is a reported diagnostic.

\paragraph{The dispatch surface.}
Fix a deviation tolerance $\delta$ and a confidence level $1-\alpha$, and let
$\tilde{q}_\tau(s,v)$, $\tau \ge 1$, denote the conformalized prediction
bounds of Eq.~\ref{eq:conformal}. These bounds carry the episode-level
marginal guarantee stated below; they are not pointwise conditional bounds at
each $(s,v)$. Set $\tilde{q}_0 \equiv 0$, since $Y_0 = 0$ identically and full
inference needs no statistical certificate. The rule is
\begin{equation}
\tau^{*}(s, v) \;=\;
\max \Bigl( \{0\} \cup \bigl\{ \tau \ge 1 \;:\; \tilde{q}_\tau(s, v) \le \delta \bigr\} \Bigr).
\label{eq:tau}
\end{equation}
The maximizing set always contains $0$, so the rule is total; when the
calibrated bounds exclude every positive depth, the step falls back to full
inference. This is the fail-closed direction made explicit in the
definition. Under (A1) the surface is bimonotone: nondecreasing in $s$,
nonincreasing in $v$. The three coarse regimes are recovered as level sets
rather than imposed: $\tau^{*} \to N$ where similarity is high and
disagreement low, $\tau^{*} = 0$ where similarity is low, and intermediate
$\tau^{*}$ where the state is covered but the neighborhood disagrees. The
last region supplies a mechanistic reading of warm start that the three-tier
rule lacked: partial denoising is retained exactly where the remaining steps
have arbitration work to do.

\paragraph{Data separation and estimation.}
The experience library is built from a dedicated prefill set
$\mathcal{D}_{\mathrm{lib}}$ and frozen before dispatch calibration. Query
episodes are then divided into three pairwise-disjoint sets---fit,
calibration, and test---none of which contributes entries to the library.
Same-trajectory retrieval is explicitly excluded. All divisions are made at
the episode level because steps within an episode are strongly dependent.

The \emph{fit} split is used to choose $k$, $W$, the depth grid, $\alpha$, and
$\delta$; fit the quantile surface; run the (A1) and
monotone-likelihood-ratio diagnostics; and perform any pilot closed-loop
selection. When $\delta$ is selected through pilot rollouts, an internal split
or cross-validation within the fit data is used. The chosen operating point
is frozen before final calibration. The \emph{calibration} split is then used
once to compute the conformal correction and is not revisited after any
methodological change. The closed-loop \emph{test} split is reserved for the
final evaluation; a full frontier may be reported, but its primary operating
point is declared in advance.
\begin{enumerate}[leftmargin=1.5em, itemsep=0.25em, topsep=0.2em]
\item \emph{Data.} For each step of the fit and calibration episodes, run
the coupled reference once and each candidate depth $\tau \ge 1$ once,
offline, recording $(s_i, v_i, Y_{i,\tau})$. This is replay of logged
states, not interaction.
\item \emph{Joint bimonotone fit.} On the fit split, estimate a single
quantile surface $\hat{q}(\tau, s, v)$, $\tau \ge 1$, by minimizing the
convex pinball (check) loss over the lattice-ordered monotone function class
of (A1). Fitting each $\tau$ separately would allow the fitted surfaces to
cross, producing depths at which a deeper skip appears safer than a
shallower one; the joint constraint removes this artifact. Order-restricted
convex optimization over a lattice partial order is solvable in polynomial
time through linear-programming or network-flow formulations.
\item \emph{Simultaneous conformal correction.} On the calibration split,
define per-episode nonconformity scores that cover all steps and all
positive depths at once,
\begin{equation}
R_e \;=\; \max_{i \in e}\; \max_{\tau \ge 1}\;
\bigl( Y_{i,\tau} - \hat{q}(\tau, s_i, v_i) \bigr),
\label{eq:R}
\end{equation}
where $i$ ranges over the steps of calibration episode $e$. Let $c$ be the
$\lceil (1-\alpha)(|\mathcal{E}|+1) \rceil$-th order statistic of
$\{R_e\}_{e \in \mathcal{E}}$ and set
\begin{equation}
\tilde{q}_\tau(s, v) \;=\; \hat{q}(\tau, s, v) + c ,
\qquad \tau \ge 1 .
\label{eq:conformal}
\end{equation}
The index must not exceed $|\mathcal{E}|$, which requires
$|\mathcal{E}| \ge (1-\alpha)/\alpha$ (at least $19$ calibration episodes
for $\alpha = 0.05$); otherwise $c = +\infty$ and the rule degenerates to
$\tau^{*} \equiv 0$, which is valid but vacuous. This is a conformalized
prediction bound with episode-level marginal coverage, rather than a
conditional guarantee at a fixed $(s,v)$. Because the correction is
simultaneous over depths and over the steps of an episode, the bound remains
valid \emph{after} the data-dependent selection of $\tau^{*}$: under
exchangeability of calibration and test episodes,
\begin{equation}
P\Bigl( \, Y_{i,\tau} \le \tilde{q}_\tau(s_i, v_i)
\ \ \text{for all steps } i \text{ and all } \tau \ge 1
\ \Bigm|\ \text{a fresh episode} \Bigr) \;\ge\; 1 - \alpha ,
\label{eq:coverage}
\end{equation}
and in particular $Y_{i,\tau^{*}} \le \delta$ holds jointly over the episode
with probability at least $1-\alpha$. Per-depth marginal calibration would
not survive this selection step; the max construction is what makes the
guarantee legitimate. The price is conservatism, which we report rather than
hide. When several tasks are pooled, Eq.~\ref{eq:coverage} is marginal
coverage over the task mixture; per-task validity requires task-wise
(Mondrian) calibration, and per-task results are reported where sample sizes
permit.
\end{enumerate}

\paragraph{Scope of the guarantee.}
Eq.~\ref{eq:coverage} applies to fresh episodes drawn from the teacher
visitation distribution represented by the calibration data. Deployment
intervenes on the process: replayed prefixes alter subsequent visitation and
may break exchangeability, reflecting the failure of composition described in
the companion note. The guaranteed quantity is also an action-deviation proxy,
not task success. We therefore study the association between $\delta$ and
closed-loop failure risk on fit data, freeze the operating point before final
calibration, and use closed-loop success on the test split as the final
criterion. The conformal result provides an offline certificate under teacher
visitation, not a guarantee for the intervened deployment distribution.

\paragraph{On the term ``training-free''.}
The estimator is a nonparametric post-hoc calibration: isotonic quantile
regression does minimize an objective over a class of monotone functions, so
it is statistical fitting in the ordinary sense. What is absent is any
policy or model fine-tuning, any gradient-based learner, and any parametric
hypothesis class; the procedure is gradient-free, order-restricted convex
optimization followed by quantile computation. We use ``training-free''
throughout in this restricted sense, consistent with the paper's standing
definition. A
learned router (E10) remains in the study as the contrast: it asks whether
gradient-based fitting buys anything beyond what calibrated order statistics
provide.

\paragraph{Relation to existing rules.}
The rule family strictly contains the incumbents: the two-threshold dispatch
is the degenerate surface that ignores $v$ and quantizes $\tau$ to three
values $\{0, \tau_{ws}, N\}$, and a fixed-depth rule is the constant surface
$\tau \equiv \tau_0$ gated by a single threshold on $s$. Containment implies
that the \emph{oracle} optimum over surfaces is no worse than the oracle
optimum over thresholds; it does not imply that a surface estimated from
finite data beats a well-calibrated threshold, since the larger class can
overfit. ``Surface $\ge$ threshold'' is therefore a hypothesis to be tested
at matched calibration budgets, not an expected theorem. In the notation of
the companion note, every surface policy is a causal scheduler over the
retrieval information set, so its replay share remains a lower bound on
$R_{\mathcal{C}(\mathcal{I}_\varphi)}(\varepsilon)$.

\paragraph{Experimental consequences.}
Evaluation proceeds in three stages. First, on one suite, we compare the
surface restricted to the depths already stored online with a calibrated
two-threshold rule under the same calibration budget. This pilot measures
whether the episode-level simultaneous correction remains useful despite its
conservatism: an overly large correction will route most states to
$\tau^{*}=0$ and eliminate any computational gain. Second, if the pilot shows
a measurable advantage, we run the full comparison with the library, teacher,
hardware, and calibration budget held fixed. The arms comprise a fixed
threshold at fixed depth, the three-tier rule, the surface, the surface with
stateful gating, a reduced-step teacher, and a learned dispatcher (E10). The
primary readout is closed-loop success against measured GPU-time per decision
step. Third, E5 reports the (A1) violation rate, the
monotone-likelihood-ratio diagnostic for the surrogate event, and the
association between deviation and closed-loop failure. All deviation data are
obtained by offline replay of logged states together with the paired initial
noise described above.

\end{document}
